\chapter{Préface}

Ce livre a une ambition simple~: un seul ouvrage cohérent de sciences de
la Terre et de l'Univers, couvrant tout de la première année d'école à la
fin de la licence, écrit pour des lecteurs du monde entier, d'abord en
anglais, et traduit ici en français.

Le style est concis et rigoureux~: des cours construits à partir de
définitions, d'exemples, de propositions, de méthodes et de schémas, avec
un raisonnement soigneux dès qu'il est accessible au niveau considéré,
suivis d'exercices choisis avec soin. Les résultats énoncés sans
démonstration complète sont explicitement marqués comme admis. Les
solutions complètes de tous les exercices sont rassemblées à la fin du
livre. Chaque nombre cité à propos de la Terre réelle, de ses océans, de
son atmosphère ou du ciel provient d'une source de référence, et chaque
courbe tracée à partir de mesures réelles est calculée d'après les données
publiées.

Le cours est réparti en cinq livres. Le premier couvre l'école primaire et
le collège, le deuxième les trois années de lycée, et les trois derniers
les trois années d'université. Les deux livres scolaires n'enseignent
jamais deux fois la même notion~: chaque idée est enseignée une seule
fois, l'année où elle trouve sa place pour la première fois, et un
chapitre ultérieur qui en a besoin s'ouvre sur un court rappel,
\emph{Tu sais déjà}, avant d'aller plus loin. Quand une notion relève de
la physique, de la biologie ou de la chimie, ce sont les livres
compagnons de la même collection qui l'enseignent, et ce livre la
rappelle avant d'adopter le point de vue du spécialiste des sciences de la
Terre. Dans les livres scolaires, chaque année est une partie, découpée en
chapitres~; chaque livre universitaire est une année, découpée en
chapitres. Plus l'année est précoce, plus le lecteur visé est jeune~:
davantage d'illustrations, des étapes détaillées plus nombreuses, et des
exercices plus progressifs.

\section*{Comment lire ce livre}

Les énoncés sont numérotés dans chaque chapitre et codés par couleur~:
les définitions en bleu, les théorèmes en rouge, les propositions et
leurs voisins en orange, et les méthodes --- courtes recettes pour les
tâches courantes --- en vert. Les exercices sont notés de ($\star$) pour
les applications directes à ($\star\star\star$) pour les problèmes plus
difficiles. Des encadrés contiennent ce qui ne fait pas partie du
raisonnement~: le rappel de ce qu'un chapitre précédent a enseigné, la
façon dont les géologues et les autres spécialistes des sciences de la
Terre travaillent sur le terrain ou au laboratoire, une page d'histoire,
des chiffres qui changent d'une année à l'autre (avec l'année pour
laquelle ils sont donnés), et des conseils de sécurité.

\section*{Sécurité}

La Terre s'apprend en la regardant~: le ciel, le temps qu'il fait, un
paysage, les pierres d'un bâtiment, la terre d'un jardin. C'est tout ce
que ce livre demande à son lecteur. Ne regarde jamais le Soleil, ni à
l'œil nu ni avec un instrument~: il rend aveugle. Les carrières, les
falaises, les grottes, les rivières en crue et les volcans sont des lieux
réservés à des professionnels formés et équipés, et ce livre décrit leur
travail sans y envoyer personne.

\section*{Contribuer}

Les sources de ce livre sont dans un dépôt git public~:
\begin{center}
\url{https://github.com/bvirrion/one-earth-and-space-sciences-book}
\end{center}
Les contributions --- nouveaux chapitres, corrections, meilleures
explications, exercices supplémentaires --- sont les bienvenues~; voir
\texttt{CONTRIBUTING.md} à la racine du dépôt.
